%%%%%%%%%%%%%%%%%%%%%%%%%%%%%%%%%%%%%%%%%%%%%%%%%%%%%%%%%%%%%%%%%%%%%%
% Lớp 11 | Xác suất | Biến cố hợp và quy tắc cộng xác suất
% Dạng 1D9?2-2: Xác định và đếm số phần tử biến cố hợp
% Số câu: 3  (sinh tự động bởi tach_cau.py; từ nay thầy sửa trực tiếp file này)
%%%%%%%%%%%%%%%%%%%%%%%%%%%%%%%%%%%%%%%%%%%%%%%%%%%%%%%%%%%%%%%%%%%%%%

%%%%%%%%%%%%%%%% TRẮC NGHIỆM 4 PHƯƠNG ÁN %%%%%%%%%%%%%%%%

% nguồn: Ok_TracNghiemLop11_Lan1
\begin{ex}%[1D9N2-2]%[ma:7edec699]
	Một tổ trong lớp 12CT có 4  học sinh nữ là Giang, Thanh, Thảo, Anh và 5 học sinh nam là Sơn, Tùng, Hoàng, Tiến, Hải. Trong giờ học, giáo viên chọn ngẫu nhiên một học sinh trong tổ đó lên bảng để kiểm tra bài. Xét các biến cố sau:\\
	H: \lq\lq Học sinh đó  là một bạn nữ\rq\rq .\\
	K: \lq\lq  Học sinh đó có tên bắt đầu bằng chữ T\rq\rq .\\
	Nêu  nội dung của biến cố hợp H hợp K.
	 \choice
		 {\lq\lq Giang\rq\rq}
		{\lq\lq Tùng, Tiến, Thảo\rq\rq} 
		{\True \lq\lq  Tùng, Tiến, Giang, Anh, Thanh, Thảo\rq\rq}
		{\lq\lq Giang,  Thanh, Thảo, Anh\rq\rq}
	 \loigiai{
	Ta có H hợp  K gồm các phần tử \lq\lq Tùng, Tiến, Giang, Anh, Thanh, Thảo\rq\rq.}
\end{ex}

% nguồn: Ok_TracNghiemLop11_Lan1
\begin{ex}%[1D9H2-2]%[ma:18c70eff]
	Gieo hai con xúc xắc cân đối và  đồng chất. Xác suất của biến cố ``Tổng số chấm xuất hiện trên hai con xúc xắc chia hết cho $5$'' là
	\choice 
		{$\dfrac{5}{36}$} 
		 {$\dfrac{1}{6}$}
		 {\True $\dfrac{7}{36}$}
		 {$\dfrac{2}{9}$}
	\loigiai{ 
	Ta có $W=\{(i; j) \mid 1 \leq i \leq 6; 1 \leq j \leq 6\}$, suy ra $n(W)=36$. 
	\\ 
	Gọi $A$ là biến cố ``Tổng số chấm xuất hiện trên hai con xúc xắc  chia hết cho 5''.\\
	Khi đó $A=\{(1; 4);(2; 3);(3; 2);(4; 1);(4; 6);(5; 5);(6; 4)\}$, suy ra $n(A)=7$.\\ 
	Do đó $P(A)=\dfrac{n(A)}{n(\Omega)}=\dfrac  7{36}$.
	 }
\end{ex}


%%%%%%%%%%%%%%%% ĐÚNG/SAI %%%%%%%%%%%%%%%%

% nguồn: Ok_dungsai_lop11_Lan1
\begin{ex}%[1D9H2-2]%[ma:eef37c0a]
 Một hộp chứa 5 viên bi xanh và 3 viên bi đỏ có cùng kích thước và khối lượng. Lấy ra ngẫu nhiên đồng thời  2 viên bi từ hộp. Gọi $A$ là biến cố \lq\lq Hai viên bi lấy ra đều có màu xanh\rq\rq, $B$ là biến cố \lq\lq Hai viên bi lấy ra đều có màu đỏ\rq\rq.
 \choiceTF 
  {\True Kết quả thuận lợi của biến cố $A$ là  $10$}
  {Kết quả thuận lợi của biến  cố $A$ là $2$}
  {\True $A \cup B$ là biến cố \lq\lq Hai viên bi lấy ra có cùng màu\rq\rq} 
  {\True Kết quả thuận lợi cho biến cố $A \cup  B$ là $13$}
 \loigiai{ 
  \begin{itemchoice}
   \itemch Số kết quả thuận lợi cho  biến cố $A$ là $\mathrm{C}_5^2=10$.
   \itemch Số  kết quả thuận lợi cho biến cố $B$ là $\mathrm{C}_3^2=3$.
   \itemch $A \cup B$ là biến cố \lq\lq Hai viên  bi lấy ra có cùng màu\rq\rq.
   \itemch Số kết quả thuận lợi cho  biến cố $A \cup B$ là $\mathrm{C}_5^2+\mathrm{C}_3^2=13$.
  \end{itemchoice}
 } 
\end{ex}
